% Faithful conversion of ages-ii-sprint-1; source wording and XML order retained.
\hypertarget{__RefHeading___Toc4787_3118176633}{}
\refstepcounter{subsection}
\subsection*{3.3.2 Sprint 1}
\label{unit:ages-ii-sprint-1}
\addcontentsline{toc}{subsection}{3.3.2 Sprint 1}

% text-p-0869
\indent Durante a Sprint 1, iniciei o trabalho dedicando-me ao estudo do Instagram Graph API com o objetivo de obter dados detalhados de perfis da plataforma, como a bio dos usuários, publicações, número de seguidores e a quantidade de posts. Esse processo exigiu uma compreensão aprofundada da API, mas obtive avanços significativos e já sabia como realizar as requisições para obter os diversos tipos de informações dos perfis.\par

% text-p-0870
\indent Além disso, conduzi uma análise minuciosa das tabelas do banco de dados da Receita Federal, categorizando todas as colunas presentes. Registrei as colunas que seriam úteis e identifiquei aquelas que não seriam utilizadas no projeto. Muitas dessas colunas continham termos codificados, então também realizei a categorização e a tradução de todos os códigos, facilitando a compreensão dos dados. Defini regras claras para a criação de novas colunas, incluindo \textit{joins} necessários para integrar os dados de maneira eficiente.\par

% text-p-0871
\indent Posteriormente, apresentei aos meus colegas toda a análise que havia realizado, detalhando minhas considerações sobre as colunas das tabelas do banco de dados. Validei junto à equipe as colunas que seriam mantidas, removidas ou questionadas aos stakeholders, além de discutir os tratamentos e enriquecimentos que seriam aplicados aos dados. Durante essa apresentação, utilizei o quadro branco para ilustrar minhas ideias e levantar questionamentos pertinentes, conduzindo um discurso ativo em frente aos meus colegas.\par

% text-p-0872
\indent Por estagiar na área de Engenharia de Dados, senti-me à vontade para assumir a liderança na parte referente ao banco de dados, dada minha experiência com extração, análise, tratamento e enriquecimento de dados. Assumi também a responsabilidade pela modelagem conceitual do nosso banco de dados. Devido à limitação de tempo, optamos por seguir diretamente para o mapeamento das entidades no Back-end, utilizando a ORM SQLAlchemy para simplificar as interações com o banco sem a necessidade de escrever SQL diretamente.\par

% text-p-0873
\indent Com o mapeamento das entidades em andamento, concretizamos os tratamentos e enriquecimentos necessários nos dados. Iniciei os tratamentos diretamente em SQL, alinhando o processo com as diretrizes estabelecidas anteriormente. Para a apresentação do projeto aos stakeholders, como ainda não tínhamos o banco de dados finalizado, me uni aos colegas AGES II e ao AGES IV Pedro Carlucci para preparar o modelo lógico relacional, que seria apresentado aos stakeholders.\par

% text-p-0874
\indent Embora não fosse comum em uma edição regular da AGES apresentar uma modelagem lógica relacional ao cliente, decidimos fazê-lo devido à importância do Banco de Dados, que era o core do projeto. A construção do banco ainda estava em andamento devido à sua complexidade, então mostramos ao cliente as modificações realizadas nas tabelas em relação ao Banco de Dados da Receita Federal, além das adições de colunas necessárias para integrar os dados das redes sociais.\par

\begin{reflection}
% text-p-0876
\indent \originaltab{}O time iniciou a sprint um tanto desnorteado. As \textit{tasks} ainda estavam em fase de definição e careciam de clareza, especialmente devido às várias incertezas do projeto, como a dificuldade de obtenção de dados das redes sociais. Devido a isso, observei muitos momentos de ociosidade por parte de alguns colegas, que se viam perdidos no andamento do trabalho. Embora eu tenha prestado auxílio pontual a alguns, deveria ter oferecido um suporte mais robusto e orientado melhor a equipe, dissipando a ‘neblina’ que pairava sobre o início da Sprint.\par

% text-p-0877
\indent \originaltab{}Ao perceber a confusão dos meus colegas e ciente da urgência das entregas, tomei a decisão de adiantar algumas tarefas por conta própria, já que eu tinha uma noção mais clara de como proceder. Convidei alguns membros da equipe para me ajudar, mas ao refletir, percebo que poderia ter sido mais proativo e envolvido um número maior de pessoas. Um suporte mais amplo teria sido benéfico para o progresso da equipe e para o aprendizado coletivo.\par

% text-p-0878
\indent \originaltab{}Apesar de um início conturbado, o time demonstrou que estava se organizando à medida em que nossas tarefas e objetivos se tornavam mais claros.\par

% text-p-0879
\indent \originaltab{}Outro aspecto que vale destacar é a responsabilidade que assumi em relação ao Banco de Dados. Desempenhei um papel relevante tanto nas decisões quanto na construção do banco e nas discussões técnicas com a equipe. Como eu estava me sentindo mais à vontade nessa área, assumi a liderança nesse ponto, orientando os colegas AGES II em questões relacionadas ao tratamento e enriquecimento de dados. Fiquei satisfeito em perceber que meu conhecimento estava agregando valor ao time, pois utilizei habilidades adquiridas fora da AGES para contribuir de maneira significativa no projeto.\par

% text-p-0880
\indent \originaltab{}Além disso, me desafiei a expor minhas ideias em frente a todos os meus colegas, utilizando o quadro branco para apresentar as propostas e conduzir as discussões. Esse momento foi crucial tanto para o meu crescimento pessoal quanto profissional, pois me permitiu sair da zona de conforto e aprimorar minha habilidade de comunicar de forma clara e assertiva em um ambiente desafiador. Ao liderar a equipe e compartilhar minhas ideias abertamente, percebi o impacto positivo que isso teve no desenvolvimento da minha autoconfiança e na minha capacidade de influenciar o time, o que reforçou minha evolução como profissional e como líder.\par

\end{reflection}
